\chapter{Bisection and the Spline}
\label{chap:bisection-and-the-spline}

% Closed question: How does a finite count of distinctions become an admissible
% notion of distance?

\begin{chapteressay}

A surveyor walking along a street paces off distance in steps. Each step
is roughly a meter. Twenty steps is approximately twenty meters. The
surveyor's distance estimate is the count of admissible steps multiplied
by the step length. The step is the unit; the count is the measurement.

The step is not infinitely refined. A step can be roughly counted; it
cannot be halved without changing the technique (a half-step is a
different gait). For ordinary distance estimation, the step is the
resolution: any distance shorter than a step is below the surveyor's
admissible precision.

If the surveyor needs finer precision, they switch to a measuring tape.
The tape is marked in millimeters. The smallest admissible distance is
one millimeter. Below this, the tape's gradations cannot distinguish.
For still finer precision, a laser interferometer is required. The
interferometer can resolve nanometers.

In every case, the distance is the count of admissible units. The unit
changes with the instrument. The count changes with the distance. The
two together produce the measurement.

This is the chapter's algorithmic perspective on distance. Distance is
a count of admissible refinements. The refinements are bisections of
the search space: each refinement halves the remaining range. The count
continues until the resolution boundary is reached --- the smallest
admissible distinction. Below the boundary, additional refinements
produce no admissible distinctions; the algorithm stops.

The chapter develops this through five algorithmic forms. The first is
bisection for root finding: the canonical example of refinement by
halving. The square root of 2 is found by 53 bisection steps to full
double-precision accuracy. The second is Galerkin finite-element
method: the analytic version of refinement, applied to differential
equations. With cubic elements, the convergence is $O(h^4)$, so a
ten-fold refinement reduces the error by 10,000. The third is
Richardson extrapolation: two refinements at different scales are
combined to cancel the leading error term, producing a stronger
estimate. The fourth is residual stopping: the algorithm terminates
when the residual falls below a calibrated threshold. The fifth is
\texttt{FINITE\_ELEPHANT} in Lean: the typeclass that formalizes the
stopping rule as a compiler-enforced certificate.

The second concrete example is the photographer focusing a manual
camera lens. The focus ring rotates continuously; the viewfinder's
image varies smoothly with the rotation. The photographer rotates the
ring to minimize the image's blur. The rotations are bisection
refinements: each correction halves the rotational uncertainty. The
photographer stops when the viewfinder's blur cannot distinguish
further refinements. The stopping is at the viewfinder's resolution.

The two examples --- bisection for $\sqrt{2}$ and lens focusing --- have
the same algorithmic structure. A continuous search space, a
bisection rule, a resolution boundary, a count of refinements until
the boundary is reached. The count is the distance. The boundary is
the calibration.

The chapter's closed question is:

\begin{center}
\emph{How does a finite count of distinctions become an admissible
notion of distance?}
\end{center}

The answer the chapter develops is: through bisection. The count of
bisection steps from a starting bracket to the resolution boundary is
the distance. The distance is finite because the resolution boundary
is finite. The distance is admissible because each bisection step is
an admissible refinement (a single binary comparison) and the stopping
rule is admissible (the resolution is part of the instrument's
calibration).

The Galerkin method generalizes bisection to higher-dimensional
problems. The basis functions are the discrete refinements; the
weak-form linear system is the algorithm; the
\texttt{FINITE\_ELEPHANT} certificate is the stopping rule.
Richardson extrapolation combines two refinements at different
scales into a stronger approximation, exploiting the convergence
structure to do better than either input alone.

The chapter closes with a connection to the next chapter. Distance
has been defined as a count of refinements. The refinements are
admissible because each one is a binary comparison certified by the
resolution. But the count itself --- which refinement, of all the
admissible ones, the algorithm should perform --- is left
unspecified. The algorithm halves the bracket; it could halve in
either direction; the choice is made by the comparison's output.
The deeper question is: when the comparison is ambiguous (two
options are equally admissible), which one does the algorithm
select? This is the motion question. Chapter 6 addresses it, and
the answer is the load-bearing chapter of the book: the selection
rule requires exact Kolmogorov complexity, which is provably
uncomputable. The compiler cannot, in general, decide which
admissible continuation to take. The universe, somehow, can.

\end{chapteressay}

\section{Bisection for Root Finding}
\label{se:bisection-for-root-finding}

To find the positive square root of 2, choose the interval $[1, 2]$ as a
starting bracket. The square root of 2 lies somewhere in this interval
because $1^2 = 1 < 2 < 4 = 2^2$. Compute the midpoint: $1.5$. The square
of the midpoint is $1.5^2 = 2.25$. Since $2.25 > 2$, the square root of
2 lies in the lower half of the interval: $[1, 1.5]$. Replace the upper
end of the bracket with the midpoint.

Repeat. New midpoint: $1.25$. Square: $1.5625$. Since $1.5625 < 2$,
the root is in the upper half: $[1.25, 1.5]$. Replace the lower end.
New midpoint: $1.375$. Square: $1.890625$. The root is in the upper
half: $[1.375, 1.5]$. And so on.

Each iteration halves the bracket. The bracket's width after $k$
iterations is $2^{-k}$. After 53 iterations, the bracket's width is
$2^{-53}$, which is the smallest interval that can be distinguished in
double-precision floating-point arithmetic. At this point, the
algorithm has produced the floating-point representation of $\sqrt{2}$
to full machine precision.

This is bisection for root finding. The algorithm requires only that
the function $f(x) = x^2 - 2$ change sign across the bracket. Given a
sign change, each iteration halves the bracket while preserving the
sign change. The bracket converges to a root.

The distance from the initial bracket to the floating-point answer is
53 bisection steps. Each step is one admissible refinement: a single
comparison decides which half of the current bracket contains the root.
The total distance is the count of steps. This is the chapter's
proposal: distance is the count of bisection refinements.

The proposal makes precise what ``close'' means algorithmically. Two
values are close when few bisection steps are required to distinguish
them. Two values are far when many steps are required. For $\sqrt{2}$
to be ``known to within $10^{-16}$,'' the bisection requires 53 steps.
For $\sqrt{2}$ to be ``known to within $0.1$,'' the bisection requires
4 steps. The number of steps is the algorithmic distance.

The connection to physical measurement is direct. A thermometer with
$0.1$ degree resolution can distinguish between $20.0$ and $20.1$. It
cannot distinguish between $20.05$ and $20.06$. The thermometer's
admissible distinctions correspond to bisection refinements of the
temperature scale. The instrument's resolution is its bisection
depth.

The bisection algorithm's correctness depends on three preconditions.
First, the function must be continuous on the bracket. Without
continuity, the sign change at the bracket does not imply a root in
the interior. Second, the function must take values of both signs at
the bracket endpoints. Without the sign change, the bracket does not
guarantee a root. Third, the function evaluation must be
\texttt{REPEATABLE}: the same input must produce the same sign every
time. Without repeatability, the bisection cannot make a deterministic
choice between the two halves.

The Lean analog of the bisection algorithm appears in the
\texttt{Polynomial} and \texttt{GalerkinProcess} typeclasses in
\texttt{Episode6.lean}. A polynomial root-finding problem is encoded
as a polynomial term and a target interval. The algorithm bisects the
interval, evaluating the polynomial at midpoints, until the
\texttt{FINITE\_ELEPHANT} stopping rule certifies that further
refinement would produce no new distinctions.

The compiler's analog of bisection is the decision-tree compilation of
pattern matching. When Lean compiles a \texttt{match} expression with
many cases, it builds a decision tree: each internal node is a
comparison; each leaf is a case body. The tree is constructed to
minimize the depth of the most common case. The compilation produces a
binary tree of comparisons, structurally identical to a bisection
search. The match expression's runtime cost is the depth of the
decision tree.

Bisection is the fundamental algorithm for refining a one-dimensional
search space. The algorithm's structure is elementary: bracket,
midpoint, comparison, half. The structure does not depend on the
specific function being evaluated or the specific space being searched.
Bisection works for square roots, for polynomial roots, for binary
search on sorted arrays, for branch decisions in compiled code, for
fare-zone determinations in transit calculations. The algorithm is the
algorithmic primitive of one-dimensional refinement.

\section{Galerkin FEM}
\label{se:galerkin-fem}

To solve the differential equation $u''(x) = -1$ on the interval $[0, 1]$
with boundary conditions $u(0) = u(1) = 0$, the analytical answer is the
parabola $u(x) = x(1-x)/2$. The maximum value is $1/8$ at $x = 1/2$.
This is a textbook problem with a closed-form solution.

The Galerkin finite element method approaches it differently. Choose a
finite-dimensional space of trial functions: piecewise cubic
polynomials with $C^2$ continuity at the knots. Place ten knots on
$[0, 1]$ at $x = 0, 0.1, 0.2, \ldots, 1.0$. The trial space has ten
basis functions (Hermite cubics). Each basis function is zero outside a
small region around its knot.

The Galerkin formulation says: find a function $u_h$ in the trial space
such that, for every test function $v_h$ in the same space, the inner
product $\int_0^1 u_h''(x) v_h(x) \, dx = -\int_0^1 v_h(x) \, dx$. The
condition reduces to a linear system: $10 \times 10$ matrix equation,
which can be solved in $O(N^3)$ time for $N$ basis functions, or
$O(N)$ for tridiagonal systems.

Solving the system produces the approximate solution $u_h$. Comparing
$u_h$ to the exact solution $u$ shows the approximation error. For $N
= 10$ cubic elements, the error is on the order of $10^{-6}$. For $N
= 100$, the error drops to $10^{-10}$. The convergence rate is $h^4$,
where $h = 1/N$ is the element size. Four orders of accuracy gained
per ten-fold refinement: the convergence is fourth-order.

This is the chapter's central convergence pattern. Each refinement
(doubling of $N$, halving of $h$) reduces the error by a factor of
$16 = 2^4$. The reduction is governed by the differential equation's
smoothness and the basis functions' polynomial degree. For cubic
elements solving a second-order ODE with smooth right-hand side, the
convergence rate is $h^4$.

The connection to distance is direct. Increasing $N$ from 10 to 100
is roughly $\log_2(10) \approx 3.3$ bisection refinements. The error
drops by $16^{3.3} \approx 5000$. The cost grows linearly: $O(N)$
operations per solve. The trade-off is: pay roughly twice as much
work to reduce the error by sixteen. For applications where the
error must be very small, the trade-off is favorable.

The Galerkin method is admissible because it produces the unique
minimum-energy approximation in the trial space. The trial space has
a specific notion of energy: the integral of the squared derivative
of the function. The Galerkin solution minimizes the error in this
energy. The minimum-energy approximation is not arbitrary; it is the
unique structure consistent with the equation and the trial space.

The Lean analog in \texttt{Episode6.lean} is the \texttt{GalerkinProcess}
typeclass:

\begin{leanexcerpt}{GalerkinProcess}
structure GalerkinProcess
    (Value : Type)
    (Carrier : CarrierProcess Value)
    [DISTINGUISHABLE Value Carrier] ... [LOAD Value Carrier]
  where
  ANSYS_process    : BASICProcess Value Carrier
  polynomial       : Polynomial
  scale_and_shift? : Polynomial $\to$ Polynomial := ...
\end{leanexcerpt}

The \texttt{ANSYS\_process} field names the underlying
\texttt{BASICProcess} that supplies the trial space. The
\texttt{polynomial} field carries the current iterate. The
\texttt{scale\_and\_shift?} method advances a polynomial by one
Galerkin step (handling constant, monomial, and factor cases). The
cascade through \texttt{LOAD} carries every prior class as a
typeclass requirement: the structure is admissible only when the
full stack has been witnessed.

\texttt{FINITE\_ELEPHANT} is the stopping rule:

\begin{leanexcerpt}{FINITE_ELEPHANT}
class FINITE_ELEPHANT
    (Value : Type)
    (Carrier : CarrierProcess Value)
    [DISTINGUISHABLE Value Carrier] ... [LOAD Value Carrier]
  where
  galerkin_process : GalerkinProcess Value Carrier
  finite?          : Polynomial $\to$ Polynomial $\to$ Prop := ...
\end{leanexcerpt}

The class binds a \texttt{GalerkinProcess} together with a binary
decision rule \texttt{finite?} on \texttt{Polynomial}s. The rule
covers the relevant transition cases: convergence to the ground
state, stationary fields (the eigenvalue hit), the Lanczos residual
check between factors, and the divergence cases. The compiler
enforces the certificate: a function that claims an answer is
admissible must provide the \texttt{FINITE\_ELEPHANT} instance.

This is the chapter's algorithmic version of the calibration
certificate. The instrument has a stated resolution. The
\texttt{FINITE\_ELEPHANT} instance certifies that the algorithm's
output is admissible at that resolution. Further refinement, if
performed, would produce a finer approximation, but the additional
distinctions would be below the instrument's resolution and
therefore not admissible as measurements.

The Galerkin method is the analytical analog of the cubic spline. The
cubic spline interpolates known data points; the Galerkin method
approximates the solution of a differential equation. Both use the
same trial space (piecewise cubic with $C^2$ continuity). Both
produce a unique minimum-energy answer. Both have $h^4$ convergence.
The cubic spline is the special case where the right-hand side is
data points; the Galerkin method is the general case where the right-
hand side is a differential equation.

The chapter's algorithmic claim is that the Galerkin method is the
algorithm forced by the chapter's question: how does a finite count
of distinctions become an admissible notion of distance? The answer:
distance is the count of refinements in the trial space, and the
trial space is the unique minimum-energy space consistent with the
admissibility constraints. The Galerkin method produces the answer;
the \texttt{FINITE\_ELEPHANT} stopping rule certifies it.

\section{Richardson Extrapolation}
\label{se:richardson-extrapolation}

The numerical derivative of $f(x) = \sin(x)$ at $x = 1$ can be
approximated by the central difference formula:
\[
f'(x) \approx \frac{f(x + h) - f(x - h)}{2h}.
\]
For $h = 0.1$: $(f(1.1) - f(0.9))/0.2 \approx 0.5377$. For $h = 0.05$:
$(f(1.05) - f(0.95))/0.1 \approx 0.5400$. The exact value is
$\cos(1) \approx 0.5403$. The error decreases as $h$ shrinks.

The error in the central difference is approximately $-h^2 f'''(x)/6
+ O(h^4)$. The leading term is proportional to $h^2$. Compute the
approximation at two values of $h$, say $h$ and $h/2$. The errors are
$-h^2 f'''(x)/6$ and $-h^2 f'''(x)/24$ respectively. Multiply the
second by 4 and subtract the first: $4 \cdot D(h/2) - D(h)$ has
error $-h^2 f'''(x)/6 - 4 \cdot h^2 f'''(x)/24 = 0$. The leading
error term cancels.

The combined formula $(4 D(h/2) - D(h)) / 3$ has error $O(h^4)$, two
orders better than either input. This is Richardson extrapolation:
combining two approximations at different resolutions to cancel the
leading error term, producing a more accurate approximation at no
additional function evaluations beyond the original two.

For $f(x) = \sin(x)$ at $x = 1$ with $h = 0.1$ and $h = 0.05$:
\[
\frac{4 \cdot 0.5400 - 0.5377}{3} = 0.5408.
\]
The exact value is $0.54030$. The Richardson value is about $5 \times
10^{-4}$ from exact. The unextrapolated $h = 0.05$ value was about
$3 \times 10^{-4}$ from exact. In this case the Richardson is
slightly worse due to the small $h$ rounding effects, but at the
generic asymptotic regime the Richardson is several orders better.

Richardson extrapolation is the algorithmic version of refinement
combination. Two approximations at different scales are combined into
a third approximation that exploits the structure of the convergence.
The structure must be known: the leading error term must depend on a
known power of $h$. For central differences, the power is 2. For the
trapezoidal rule, also 2. For Simpson's rule, the power is 4. For
each method, Richardson is the canonical accelerator.

The connection to the chapter's question is that Richardson
extrapolation produces a distance estimate that exploits the
algorithm's known convergence rate. The estimate's accuracy is not
just the resolution at the finer scale; it is the projected
accuracy from extrapolating the error toward $h = 0$. The
extrapolation projects a refinement that the algorithm did not
perform, but that the error structure permits.

The Lean analog is the typeclass that combines two
\texttt{GalerkinProcess} instances at different resolutions. Given
both, the algorithm can compute a Richardson estimate. The
\texttt{FINITE\_ELEPHANT} certificate for the Richardson result is
stronger than the certificate for either input: it certifies
admissibility at a resolution finer than either input could provide
on its own.

The chapter's algorithmic claim is that Richardson extrapolation is
the algorithmic primitive for combining refinements. Two
refinements alone are just two refinements; their combination, via
Richardson, exploits the convergence structure to produce a third
refinement at a stronger resolution. The third refinement is not a
new measurement; it is an algebraic combination of the two
measurements that the algorithm did perform. The combination is
admissible because the convergence structure is part of the
algorithm's specification.

Richardson extrapolation works in higher dimensions too. For a
two-dimensional grid with $h$ being the grid spacing, the error
structure depends on $h^2$ in many algorithms (e.g., the finite
difference Laplacian on a uniform grid). Refining the grid by half
and applying Richardson produces a fourth-order accurate Laplacian
from a second-order accurate one. The technique is also used in
numerical integration (Romberg integration), in ODE solvers
(Bulirsch-Stoer method), and in machine learning model training
(model averaging across step sizes).

In every case, the algorithmic content is the same. Two
approximations at different resolutions, combined to cancel the
leading error term, produce a higher-order approximation. The
combination is admissible because the convergence structure is
known and the cancellation is exact.

\section{Residual Stopping}
\label{se:residual-stopping}

The Conjugate Gradient method solves the linear system $Ax = b$, where
$A$ is a symmetric positive-definite matrix, by iteratively improving
an estimate $x_k$. At each iteration, the algorithm computes the
residual $r_k = b - A x_k$, the difference between the current
estimate's effect and the desired right-hand side. If $r_k$ is small,
$x_k$ is close to the true solution.

The iteration continues until the residual's norm $\|r_k\|$ falls below
a threshold $\varepsilon$. The threshold is the stopping criterion. The
algorithm has reached an admissible approximation when the residual no
longer carries enough information to refine the solution further. Below
the threshold, additional iterations would adjust the solution by
amounts smaller than the threshold, producing changes that are
invisible at the calibrated precision.

This is the algorithmic form of resolution-based stopping. The
algorithm runs as long as the residual is producing distinctions at
the current scale; it stops when the residual falls below the scale's
resolution. The stopping condition is not a fixed number of iterations;
it is a relationship between the residual and the resolution.

The choice of $\varepsilon$ is part of the calibration. A scientific
application that requires four significant figures sets
$\varepsilon = 10^{-4}$; an engineering application that requires only
two sets $\varepsilon = 10^{-2}$. The threshold determines the
admissible accuracy of the output. Below the threshold, the algorithm
admits the answer as final. Above the threshold, the algorithm
continues iterating.

The stopping criterion's algorithmic content is that it terminates
the refinement process at the resolution boundary. The chapter's
question --- how does a finite count of distinctions become an
admissible notion of distance? --- has its answer in part here. The
distance is the count of iterations the algorithm performs before the
residual falls below the threshold. The count depends on the
threshold, on the matrix's properties (condition number, spectral
distribution), and on the initial guess. The count is the
algorithm's measurement of the problem's distance.

For well-conditioned matrices (condition number close to 1), the
Conjugate Gradient converges quickly: the count is small. For
ill-conditioned matrices (condition number much greater than 1), the
count is large. The count is the chapter's distance, measured in
units of admissible refinements.

The Lean analog is the \texttt{FINITE\_ELEPHANT} typeclass mentioned
in the Galerkin section. The class certifies that the algorithm's
current state is admissible at the resolution. The certification is
the algorithmic version of the residual-falls-below-threshold check.
The compiler enforces the certificate: a function that asserts an
answer is admissible at a stated resolution must provide a
\texttt{FINITE\_ELEPHANT} instance for that resolution.

Residual stopping appears in many iterative algorithms. Newton's
method for root finding stops when the function value at the current
estimate is below a threshold. Gradient descent for optimization stops
when the gradient norm is below a threshold. Iterative refinement of a
linear solve stops when the residual is below a threshold. Expectation-
maximization in statistical estimation stops when the change in
log-likelihood is below a threshold. In each case, the algorithm
terminates when further iteration would produce no admissible
refinement.

The chapter's claim is that residual stopping is the algorithmic
form of resolution-bounded refinement. The stopping condition is the
calibration boundary: above the boundary, the algorithm has admissible
work to do; below the boundary, the algorithm has produced an
admissible answer. The boundary is the resolution. The boundary is
part of the algorithm's specification. Without the boundary, the
algorithm could iterate forever, producing increasingly tiny
adjustments that are invisible at the calibrated precision.

The compiler analog of residual stopping is the elaboration
heartbeat limit. Lean's \texttt{set\_option maxHeartbeats 400000}
caps the elaborator's effort at 400,000 heartbeats per definition.
When the limit is reached, the elaborator stops; the definition
fails to elaborate. The heartbeat limit is the resolution boundary
of the compiler: definitions whose elaboration cost exceeds the
limit are not admissible. The user must either simplify the
definition or raise the limit.

The heartbeat limit is the compiler's calibration. Different
projects choose different limits based on the complexity of their
proofs. Mathlib uses a relatively high limit because its proofs are
complex; smaller projects use lower limits. The limit is the
admissible resolution of the elaborator. Definitions that exceed
the limit are too coarse for the calibration; the user must refine
them.

This connects the chapter's algorithmic notion of distance to the
compiler's specific implementation. Distance is the count of
admissible refinements. The compiler's distance is the count of
heartbeats. The chapter's algorithmic claim is that this is the
correct generalization: distance is whatever count the calibration
specifies as admissible. The specific units (bisection steps,
Galerkin refinements, Conjugate Gradient iterations, compiler
heartbeats) vary, but the structure is the same.

\section{FINITE\_ELEPHANT in Lean}
\label{se:finite-elephant-in-lean}

The typeclass \texttt{FINITE\_ELEPHANT} in \texttt{Episode6.lean} is the
formalization of the chapter's stopping rule. The class certifies that
a finite element approximation is admissible at its stated resolution:
further refinement would produce no new distinctions.

The class has the shape:

\begin{leanexcerpt}{FINITE_ELEPHANT}
class FINITE_ELEPHANT
    (Value : Type)
    (Carrier : CarrierProcess Value)
    [DISTINGUISHABLE Value Carrier] ... [LOAD Value Carrier]
  where
  galerkin_process : GalerkinProcess Value Carrier
  finite?          : Polynomial $\to$ Polynomial $\to$ Prop := ...
\end{leanexcerpt}

The first field, \texttt{galerkin\_process}, names the underlying
Galerkin iteration. The second field, \texttt{finite?}, is the
decision rule that classifies which transitions are admissibly
finite. The cascade obligation runs the entire length of the class
stack; the rule covers the transition cases the chapter named.

To provide a \texttt{FINITE\_ELEPHANT} instance for a specific
\texttt{GalerkinProcess}, the user must witness every class in the
cascade and supply (or accept the default for) the decision rule.
The witnesses typically come from convergence theorems: for cubic
elements solving a smooth ODE, the error is $O(h^4)$; with
$h = 0.01$, the error is $10^{-8}$; the rule classifies the
resulting transitions as finite.

The chapter's algorithmic claim is that \texttt{FINITE\_ELEPHANT} is
the compiler-enforced version of the stopping rule. Without the
certificate, the compiler will not admit the approximation as final.
With the certificate, the approximation is admitted, and downstream
computations can use it as a record.

The compiler's enforcement is at elaboration time. If a function
asserts that an approximation is admissible at resolution
$\varepsilon$ without providing a \texttt{FINITE\_ELEPHANT} instance,
the compiler refuses to admit the function. If the instance is
provided but the proof of admissibility is incorrect, the compiler
checks the proof and produces an error.

This is the chapter's resolution-based stopping in formal Lean code.
The algorithm's distance is the count of refinements; the count
terminates at the resolution boundary; the \texttt{FINITE\_ELEPHANT}
certificate marks the termination as admissible. The certificate is
the algorithm's commitment that its output is final at the stated
resolution.

The accompanying typeclasses extend the framework. \texttt{Spline}
is the type of cubic spline approximations. \texttt{Spline.le} is the
partial order on splines: one spline refines another when its
resolution is finer (smaller error tolerance). The order makes
splines composable: a coarser spline can be replaced by a finer one
that satisfies stricter admissibility. The replacement does not lose
information; it adds distinctions.

\texttt{MAGNITUDE}, \texttt{SCALED}, and \texttt{LOAD} are the
typeclasses for the algebraic operations on splines:
\texttt{MAGNITUDE} for the size of a spline value; \texttt{SCALED}
for multiplication by a scalar; \texttt{LOAD} for the right-hand side
of the differential equation being approximated. Together they
formalize the linear algebra that makes the Galerkin method work.

The chapter closes with this observation. Distance, in the
algorithmic sense, is a count of admissible refinements. The count
is bounded by the resolution. The resolution is part of the
algorithm's calibration. The \texttt{FINITE\_ELEPHANT} typeclass is
the formal commitment to the resolution. Together they constitute
the algorithm's promise: the output is admissible at the stated
resolution, and the algorithm has terminated because further
refinement would produce no admissible distinctions.

The chapter's argument is that this discipline is forced by the
chapter's question. To convert a finite count of distinctions into
an admissible notion of distance, the algorithm must declare its
resolution explicitly. Without the declaration, distance is
underspecified; with the declaration, distance is the count of
refinements until the resolution is met. The \texttt{FINITE\_ELEPHANT}
typeclass is the formal version of the declaration. Programs that
provide it are certified; programs that do not are not yet
algorithmically complete.

\begin{bridge}

The chapter's question was how a finite count of distinctions becomes
an admissible notion of distance.

The answer is: through bisection, with the count of refinements
serving as the distance. Bisection halves the search space at each
step; the count to reach the resolution boundary is the distance.
Galerkin FEM applies the same structure to differential equations:
the count of basis functions refines the approximation, with $h^4$
convergence for cubic elements. Richardson extrapolation combines
two refinements to cancel the leading error and produce a stronger
approximation. Residual stopping terminates the refinement when
further iteration would produce no admissible distinctions.
\texttt{FINITE\_ELEPHANT} formalizes the stopping rule as a compiler-
enforced typeclass.

What remains is motion: how does the system select one admissible
continuation from among the many that the refinement allows? The
next chapter is the load-bearing chapter of the book. It contains
the noncomputability argument: the selection rule requires exact
Kolmogorov complexity, which is provably uncomputable. The
algorithmic vocabulary built in the first five chapters is the
preparation; Chapter 6 is the argument.

\end{bridge}

\begin{coda}{Focusing a Camera Lens}

The photographer is standing on a rooftop at twilight. The subject is
a building three blocks away: a clock tower with an illuminated face.
The camera is on a tripod. The lens is a 200-millimeter telephoto,
manual focus. The photographer's eye is at the viewfinder.

The focus ring is a smooth metal cylinder marked with distance
gradations: 5 meters, 10 meters, 20 meters, infinity. Between the
gradations, the ring rotates continuously. The photographer's job
is to find the position of the ring at which the clock tower is in
sharpest focus.

He starts with the ring at infinity. Through the viewfinder, the
clock face is blurry. He rotates the ring counter-clockwise. The
blur decreases, then begins to increase. He has overshot. He
rotates back clockwise. The blur decreases again. He has
overshot in the other direction.

This is bisection. The bracket is the range of focus settings
between the two overshoots. The midpoint is the photographer's
next estimate. Each rotation halves the bracket.

For the first pass, the bracket is about thirty degrees of ring
rotation, corresponding to a distance range of two to three
meters (the tower is three blocks away, but the lens's focus
range is large at the long-focus end). The photographer makes
three or four corrections, each one halving the rotational
bracket. After the corrections, the bracket is about two degrees
of rotation.

At this point, he refines the technique. He inspects the
viewfinder more carefully. The clock face has a long minute
hand and a short hour hand. The minute hand is the finer detail;
he focuses on it. A small turn of the ring sharpens the minute
hand, then begins to soften it. He has reached the bisection's
fine scale: rotations smaller than a degree.

The lens's optics are such that the image on the camera's sensor
varies smoothly with the focus ring's position. The variation is
the convergence: small changes in the ring produce small changes
in the image. The smoothness is what makes bisection work. If
the variation were discontinuous, the photographer would not be
able to converge by halving the bracket.

At ten degrees of bracket, the photographer can see the
distinction between sharp and slightly-soft. At one degree, he
can still see it, but the distinction is at the limit of his
visual discrimination. At a tenth of a degree, the distinction
is invisible to him through the viewfinder. The clock face
looks identical at the center of the bracket and at the
endpoints.

This is the resolution boundary. The viewfinder is the
calibration instrument. Below the resolution, additional
refinement of the focus ring produces no admissible distinction.
The lens may technically be slightly out of focus by a tenth of
a degree, but the photographer cannot tell, and the final
photograph will not reveal a difference (the camera's sensor
has its own resolution boundary, somewhat coarser than the
photographer's eye).

The photographer stops. The focus is locked. The distance from
the initial setting (infinity) to the final setting is the
count of bisection steps: roughly twelve corrections, each
halving the bracket. Each correction is an admissible
refinement. The total distance is twelve refinements.

If a more precise focus is required --- say, for a macro shot
of a small object --- the photographer would need additional
tools. A focus rail allows millimeter-level adjustments of the
camera's distance from the subject. A focus loupe magnifies the
viewfinder, exposing finer distinctions. With these tools, the
bisection continues to a finer scale: the resolution boundary
moves. But there is always a resolution: the camera's sensor
has finite pixels, the lens has a finite resolution power, the
photographer has finite visual acuity. Below the boundary, the
image cannot be improved by further refinement.

This is the chapter's lesson. Distance is the count of
bisection steps to the resolution boundary. The boundary is
real, not arbitrary. Different instruments have different
boundaries; the same instrument under different conditions has
different boundaries (the photographer can focus more precisely
in bright light than in dim light). The boundary is part of
the instrument's specification. The
\texttt{FINITE\_ELEPHANT} certificate is the formal version of
``focus is now locked at the resolution''.

The photographer presses the shutter. The image is captured.
The final image's resolution is the joint resolution of the
lens, the sensor, and the focus accuracy: each one contributes
to the total. If any one is worse than the others, it
dominates the result. The lens may be a fine optic; the sensor
may have many megapixels; if the focus is two degrees off, the
image is soft and no amount of post-processing recovers the
detail.

The bisection algorithm of focusing terminates because the
viewfinder cannot show finer distinctions. The Galerkin
analog terminates because the residual is below the threshold.
The Richardson analog terminates when the extrapolation
agrees with the unextrapolated value. The Conjugate Gradient
analog terminates when the residual is below the
$\varepsilon$ threshold. In every case, the algorithm
terminates at the calibrated boundary of admissible
refinement.

On the rooftop, the photographer takes three or four more
exposures of the clock tower at the same focus. The exposures
differ in their light conditions; the focus does not change.
The bracket of light conditions is being explored; the focus
bracket has been resolved. The instrument is at its
resolution; the next measurement (light) has its own
resolution. The two are decoupled: the focus is a
\texttt{FINITE\_ELEPHANT} on the focus axis, and the
exposures iterate the light axis. The two axes have their
own bisections, their own resolutions, their own
\texttt{FINITE\_ELEPHANT} certificates.

The photographer packs up at full dark. The day's photographs
are stored. Tomorrow he will return to the same rooftop, with
the same camera, at the same time. The focus ring will be at
the same setting --- the calibration is repeatable. The
exposures will differ as the clouds and lights change. The
bisection at each axis will repeat. The
\texttt{REPEATABLE} property of the instrument allows the
algorithm to be applied repeatedly, producing comparable
results from session to session. The chapter's
algorithmic forms compose: \texttt{REPEATABLE} comparison
underpins the \texttt{FINITE\_ELEPHANT} stopping rule, and
the \texttt{FINITE\_ELEPHANT} certifies the resolution of
the admissible refinement.

\end{coda}
